Un noyau de nombre de masse $A = \v{A}$, initialement au repos, subit une désintégration alpha. La particule alpha (masse \qty{6.645e-27}{kg}) est émise avec une énergie cinétique de \qty{\v{K}}{MeV}. On utilise la mécanique classique et on approxime la masse de chaque noyau par son nombre de masse.

\begin{enumerate}[label=\alph*)]
    \item Déterminez la vitesse de la particule alpha.
    \item En appliquant la conservation de la quantité de mouvement, déterminez l'énergie cinétique de recul du noyau fils.
    \item Déterminez l'énergie totale libérée par la désintégration (valeur $Q$).
\end{enumerate}
